\section{Limitations}

The evidential scope is intentionally narrower than the general world-model
problem. First, all trajectories come from five analytically specified,
low-dimensional controlled systems. This gives exact split, corruption, and
parameter-shift provenance, but it does not reproduce the visual ambiguity,
contacts, unmodeled objects, sensor synchronization failures, or nonstationary
actuation encountered by embodied agents. The state is directly observed
before synthetic corruption, future controls are supplied, and no policy is
trained. Prediction quality here therefore cannot be converted into a claim
about control return, planning safety, or real-robot transfer.

Second, every model receives the correct differential-equation family with
imperfect parameters. Structural misspecification, hidden state variables, and
informative sensor dropout may change the balance between nominal propagation
and measurement correction. The gain 0.5 is deliberately fixed rather than
estimated from uncertainty; the gain sweep measures sensitivity on this
benchmark but does not establish filter optimality. The Gaussian diagonal
forecast likewise cannot represent multimodality or full cross-state
covariance, and the bounded one-step residual does not guarantee stable
multi-step trajectories.

The two corrupt evaluation regimes bundle Gaussian observation noise and
missingness, and the OOD-corrupt regime additionally bundles parameter shift.
There are no noise-only, missingness-only, or clean-OOD factorial cells.
Accordingly, regime differences cannot identify which perturbation caused an
effect or establish the complementary update as its causal mechanism.

Third, inference has only eight independent seed-level training/data
realizations. The exact sign-flip test is attainable for the prespecified
five-comparison families, but it has coarse resolution and assumes symmetry of
seed effects. Trajectory resampling improves precision for the fixed benchmark
without creating additional independent training replicates. System--regime
contrasts are consequently descriptive, and all generalization claims remain
limited to the frozen five-system grid. Larger preregistered seed studies,
structurally wrong physics, state-dependent missingness, visual encoders, and
hardware-in-the-loop evaluation are necessary before deployment claims would
be justified.
